\section{Introduction}

\subsection{Background and Motivation}
The integration of renewable energy sources into the power grid has driven the development of Virtual Power Plants (VPPs) to aggregate and manage distributed energy resources. Accurate forecasting of VPP dynamics---specifically electrical load, photovoltaic (PV) generation, and wind power generation---is critical for grid stability, economic dispatch, and operational risk management \cite{Ma2025,Cao2023,Qiu2024,Moreno2020,Wu2024}. Unlike single-variable forecasting, cooperative VPP management requires simultaneous, coherent predictions across these interrelated domains, where the physical constraints of each energy modality must be rigorously respected. For instance, PV generation is intrinsically bound by solar irradiance and cannot produce power during nighttime, while wind generation depends on a non-linear turbine power curve driven by localized wind speed. 

Historically, early approaches to VPP forecasting relied on statistical extrapolation. As distributed resources increased, recurrent neural networks dominated sequential modeling but struggled with long-range dependencies. Subsequently, Transformer-based frameworks emerged as the dominant paradigm, capturing varied temporal dynamics across extended sequences \cite{Zhou2021informer,Wu2021autoformer,Nie2023patchtst,Zeng2023}. Within the VPP domain, multivariable Transformer adaptations have demonstrated competitive accuracy \cite{Qiu2024,Ashraf2025}. Reliable dispatch decisions hinge not only on minimizing raw statistical forecasting errors but also on ensuring that predictive models fundamentally comply with the immutable physical laws governing these grids.

\subsection{Limitations of Existing Data-Driven Methods}
Recent advancements in deep learning have introduced powerful time-series forecasting architectures, yet their application to VPP scheduling presents fundamental challenges across three distinct dimensions \cite{Vle2025,Arabzadeh2025}:

\textbf{1) Mathematical Black-Box Oblivion:} While deep neural networks excel at minimizing statistical error metrics, they operate as black-box predictors that remain oblivious to the underlying physical principles governing power systems. Consequently, pure data-driven approaches optimize exclusively for statistical error minimization, remaining unconstrained by thermodynamic laws and frequently producing physically impossible predictions---such as non-zero solar power during the night or negative wind generation. These deficits severely undermine operational safety and operator trust.

\textbf{2) Fragility of Post-Hoc Physics Interventions:} To address physical non-compliance, Physics-Informed Neural Networks (PINNs) \cite{Raissi2019} incorporate physical differential equations as soft constraints. However, VPPs operate in highly stochastic environments where a device's active status depends on volatile exogenous variables rather than on fixed deterministic rules \cite{Li2026}. Fixed physical parameters and static penalty terms fail to adapt dynamically to transient wind or solar phenomena \cite{Asinyo2025}. Alternatively, applying heuristic post-processing---such as hard zero-clipping at physical boundaries---can retroactively suppress the violation rate to zero but strictly functions as a superficial patch. Heuristic clipping inherently disrupts end-to-end optimization by introducing non-differentiable dead zones, severing the natural continuous manifold of the underlying predictions without addressing the model's structural ignorance.

\textbf{3) Lack of Structural Traceability:} As forecasting models grow increasingly complex, structural interpretability represents a critical safety requirement \cite{Vle2025}. Prior interpretability efforts relied on post-hoc attribution or internal attention weight visualization, which do not provide deterministic physical causal attribution \cite{Lovo2025}. In operational contexts, grid operators require models that explicitly decode how specific, measurable environmental variables govern device activation bounds. A persistent gap remains in causal gating formulations capable of dynamically linking network activations to observable physical state transitions \cite{Song2025,Chen2024}.

\subsection{Proposed Approach and Main Contributions}
Bridging the gap between the high predictive capacity of deep learning and the rigorous physical compliance required by coupled power systems remains a critical challenge. To this end, this paper proposes \textbf{PhysFormer}, a physics-guided causal Transformer designed to natively integrate differentiable physical priors upstream of the forecasting output, anchoring the neural network to intrinsic structural compliance.

The proposed architecture fundamentally breaks the dichotomy between predictive acuity and structural rigor. The primary contributions of this paper are twofold:

\begin{enumerate}
    \item A \textbf{Bounded Physical Act-Residual (BPAR)} architecture is proposed. By fusing parallel Transformer and ExplicitPhysicalMapping streams via Softplus lower-bound clamping and physics activity masks, this mechanism intrinsically limits the neural residual search space. It mathematically guarantees that forecasting output never breaches normalized zero bounds, eradicating the boundary violation rate (BVR=0.00\%) without resorting to heuristic post-processing.
    
    \item An orthogonal \textbf{Physics-Guided Causal Coupling (PGCC)} mechanism is introduced. By explicitly regularizing the cross-attention irradiance threshold learning, PGCC successfully imprints a highly correlated physical relationship onto the temporal gating variable. This structural alignment separates attribution transparency from error minimization capacity, conferring quantifiable structural transparency and extreme resilience to out-of-distribution (OOD) sensor noise.
\end{enumerate}

To ensure optimal learning of these components, we additionally define a \textbf{Physics-Constrained Curriculum Training} strategy that seamlessly optimizes numerical regression alongside rigid engineering bounds to ensure absolute multi-objective saturation. Furthermore, we redefine the \textbf{Ramp Violation Metric (RVM)} within a \textit{safety-first} framework, demonstrating that PhysFormer successfully bounds minor structural corrections at boundary limits to physically negligible magnitudes, thereby prioritizing absolute safety over blind continuity.


